\section{Boundary structure and family theorems}
\label{app:boundary-overview}

The following statements organize the independent boundary and family
results proved in the subsequent appendices.  They retain the full
primary and specialization theory while separating it from the
logical spine of the inverse theorem.

\subsection{Sharp global laws}

The general exponent has a sharper geometric form on the open
\(G_4^\circ\).  Theorem~\ref{thm:sharp-global-jacobian-depth} proves
\[
 \widehat{\mathcal N}_R^{\,5}=0,
 \qquad \widehat{\mathcal N}_R^{\,4}\ne0
 \ \text{at every point of projection corank three or four}.
\]
The reason is global rather than semicontinuous: the determinant-four
identity in the full matrix coordinate ring restricts to every Schur
chart.  The needed coordinate-ring Pieri projection is calculated with
the Fischer form; its coefficient is \(1/35\), not an unspecified
straightening coefficient.  The lower bound at corank three is the
Schur-symbol exclusion proved below.  The coefficient-space tables
are retained in full, but Proposition~\ref{prop:smooth-web-rank-admissibility}
determines necessary rank conditions for their occurrence on the
smooth, basepoint-free geometric open.

There is also a global geometric description beyond determinant depth.
Let \(Z_R\) be the image of the incidence vector bundle
\(\operatorname{Hom}(V,\mathcal H)\) over the Jacobian K3 surface,
and let \(D_2\) be the rank-at-most-two matrix variety.
Theorem~\ref{thm:global-minimal-supports} identifies these as the
unique minimal supports of the first and second layers.  The incidence
resolution has fibres given by linear sections of the same K3:
four contact points, a plane quartic of genus three, and the entire
surface.  Theorem~\ref{thm:cross-corank-incidence} realizes this
sequence in one two-parameter family and computes its seven reduced
components.  Proposition~\ref{prop:cross-corank-graded-fibres}
computes the intrinsic graded residue-field algebras across the same
corank transitions.  These statements do not assert a complete list
of all embedded transverse primes of the third and fourth layers.

The relative theorem is proved by primary modules, rather than by an
inference from open-containment constructibility.  A diagonal
injection into filtered primary quotients excludes every untracked
associated point; associated-element injections give the opposite
inclusion.  A regular element on the quotient by a coherent torsion
module prevents new packet torsion after specialization.  Both kinds
of injections survive by flattening their actual cokernels.
In the proper quotient-induced family, the classical
Bruns--Vasconcelos identity \cite{BrunsVasconcelos} further gives
\(dJ=(d)\cap P^{a+1}\), an exact exponent and explicit embedded
lengths (Theorem~\ref{thm:quotient-induced-primary-law}).  The
classical minor identity itself is not a novelty claim of this paper.

The argument has two scales.  The first is universal: determinant
completion bounds the principal nilpotence exponent in every symmetric
power.  The second is relative: once such an exponent is bounded, the
entire colon tower is finite and can be stratified simultaneously.
Combining the two produces a uniform theorem on the parameter space of
surjections \(\gamma:\Sym^rV\twoheadrightarrow S\): after finite
stratification, all residual colons, geometric associated-point packets,
embedded multiplicities, and multiplication ranks are compatible with
arbitrary base change.
The bounded principal-colon theorem proved below makes kernel formation
commute with arbitrary base change before any associated-point
argument is invoked.  Actual primary quotients then exhaust the geometric associated set.
Orbit-ideal descent and a power-certified construction on the
original base produce the closed packets; a relative torsion module
controls their embedded multiplicities.
For the \((1,4,6)\) problem this mechanism is supplemented by exact
rational-function-field calculations on transverse stabilizer slices
and by explicit boundary degenerations.



\subsection{The intrinsic graded boundary packet}

Write \(\widehat\NN_R\) for the nilradical of
\(\OO_{\widehat D_R}\).  The schemes
\[
 \ZZ_j(\widehat D_R)=
 V\!\left(\Ann(\widehat\NN_R^{j})\right)
\]
and the products among
\(\widehat\NN_R^j/\widehat\NN_R^{j+1}\) are invariant under
every abstract isomorphism of \(\widehat D_R\).  On a Schur chart
the reduced equation is \(d=\det T\) and
\[
 \mathcal I_{\widehat D_R}=dJ.
\]
The residual-colon theorem of Section~\ref{sec:stratified-nilpotent}
identifies
\begin{equation}\label{eq:intro-graded-packet}
 \widehat\NN_R^j/\widehat\NN_R^{j+1}
 \simeq
 \mathcal L^{\otimes j}\otimes\OO_{W_j},\qquad
 W_j=V\bigl((d)+(J:d^{j-1})\bigr).
\end{equation}
The coordinate expression changes from chart to chart; the modules
and schemes in \eqref{eq:intro-graded-packet} do not.

\begin{theorem}[Boundary structure]
\label{thm:main-depth}
For \(R\in G_4^\circ\) the following hold.
\begin{enumerate}
\item On every projection-rank chart,
\[
 d^5\in J,\qquad \widehat\NN_R^6=0.
\]
This is the \((e,r)=(4,2)\) case of the universal determinant
completion theorem, whose exponent
\(\binom{e+r-1}{r-1}\) is sharp in the universal class.

\item The multiplication of the positive graded pieces is the
canonical quotient multiplication
\[
 \OO_{W_i}\otimes\OO_{W_j}\longrightarrow\OO_{W_{i+j}}
\]
with the expected conormal-line twist.  No additional scalar
multiplication constants occur.

\item If the pure block \(A_H\) is injective at projection corank
two, the mixed block has nine orbit types.  On the generic
coefficient stratum of every orbit, the complete nonzero graded
supports through \(W_3\), their Hilbert functions, radicals,
associated supports, and vertex-primary corrections are given in
Theorems~\ref{thm:corank-two-mixed-kernel-atlas} and
\ref{thm:generic-primary-signatures}.  Genericity is proved over
rational function fields on transverse stabilizer slices.

\item If \(A_H\) drops rank at projection corank two, every
generic rank type \((a,b)\) has an explicit primary signature.
Besides the Segre quadric, the only positive-dimensional generic
supports that occur are reduced divisors of two or four lines in one
ruling.  The remaining layers are vertex-primary; in the
\((a,b)=(1,3)\) row a length-three embedded vertex accompanies
the two ruling lines.

\item At projection corank three, \(d^3\notin J\) and
\(d^4\in J\); the index is exactly five.  The upper bound is the
restriction of the full-matrix identity, as proved in
Theorem~\ref{thm:sharp-global-jacobian-depth}.  The lower bound is
the functorial Schur-symbol exclusion.

\item At projection corank four,
\[
 d^3\notin J,\qquad d^4\in J.
\]
The second membership is proved from the actual Cauchy--Pieri
multiplication in \(\C[\End(V)]\): the product of the
\((4,4,4,0)\)- and \((4)\)-components has nonzero
\((4,4,4,4)\)-projection.  Since that line is spanned by
\(d^4\), the local nilpotency index is exactly five.
\end{enumerate}
\end{theorem}

Thus the phrase \emph{boundary atlas} refers to intrinsic graded
data, not to the chosen coefficient matrices.  The Macaulay matrices
form a finite presentation atlas.  Any two such presentation atlases
have a common refinement and compute the same intrinsic primary
signature.


\subsection{Primary specialization}

The preceding finite calculations are promoted to family statements
in Section~\ref{sec:relative-primary-specialization}.  The key
homological step is explicit.  If \(B=P_S/J\),
\(C_j=B/d^{j-1}B\), and \(I_j=d^{j-1}B\), we first flatten
\(B\) and \(C_j\).  The exact sequence
\(0\to I_j\to B\to C_j\to0\) then makes \(I_j\) flat, and
\(0\to K_j\to P_S\to I_j\to0\) remains exact after arbitrary
base change.  Only after the equality
\((K_j)_s=(J_s:d_s^{j-1})\) has been established do we apply
relative assassins.  Finite primary models and the two opposite associated-point inclusions
give exact geometric support packets.  A regular element on the quotient
by a coherent torsion module excludes new torsion after specialization,
and a finite power filtration computes the local lengths.  The
closed packets descend as ideals with an explicit cocycle in
characteristic zero.  In arbitrary characteristic their ideals are
constructed downstairs first and their geometric components are
certified by finite power inclusions on a splitting cover.  The
relative torsion module is also constructed downstairs.  Noetherian
induction then handles the complement.

The same section writes explicit one-parameter degenerations of the
mixed block.  For example an off-Segre kernel point specializes to a
Segre point; a secant kernel line specializes to a tangent or either
ruling line; rank-two kernels specialize to the two rank-one types;
and both rank-one types specialize to the zero mixed map.  The
graded schemes in the two fibres are computed, so the specialization
records the birth of vertex-primary and ruling components rather than
merely changing orbit labels.


\subsection{The simple-contact locus}

The general atlas contains the earlier primary theorem as its split
form.  We retain the following statement because it is the geometric
entry point to the collision analysis.

\begin{theorem}[Corank-two primary structure on the simple-contact open]
\label{thm:main-coranktwo}
There is a nonempty invariant open
\(G_4^\dagger\subset G_4^\circ\) such that, for each
\(R\in G_4^\dagger\), a dense open set of rank-two images
\(H\subset V\) has, after an etale extension and smooth auxiliary
coordinates,
\begin{equation}\label{eq:main-coranktwo}
 \mathcal I_{\widehat D_R}
 =
 (\delta)\cap P_1^2\cap P_2^2\cap P_3^2\cap P_4^2
 \cap\mathfrak m^5.
\end{equation}
The rank-two stratum is an additional associated support and
\(\widehat\NN_R^3=0\ne\widehat\NN_R^2\).
\end{theorem}

Repeated roots, the containment locus, the decomposable-kernel wall,
and all remaining mixed-kernel orbits are treated without splitting
the contact quartic.


\subsection{Variation and finite ambiguity}

\begin{theorem}[Variation and finite ambiguity]
\label{thm:main-moduli}
For \(e=4\), the image of
\[
 R\longmapsto(Y_R,\OO_{Y_R}(1))
\]
in quartic K3 moduli has dimension nine.  On a nonempty invariant
open the induced map from projective equivalence classes of relation
spaces to its image is generically finite.
\end{theorem}

The nine-dimensional count is classical Reye geometry, not a novelty
claim.  The polarized K3 alone is not identified with the original web.
Theorem~\ref{thm:generic-intrinsic-web} instead uses the deepest
nonreduced normal cone to recover every smooth Jacobian web, and
Corollary~\ref{cor:separate-torelli-packet} separates the finite
packet.  The degree of the map defined by the K3 alone is not asserted
to be one.  The sharp exponent, the two minimal supports, and the
explicit cross-corank geometry of
Section~\ref{sec:sharp-global-laws} remain independent global
consequences of the primary-boundary analysis.


